\documentclass[11pt,a4paper]{article}

\usepackage[utf8]{inputenc}
\usepackage[T1]{fontenc}
\usepackage{lmodern}
\usepackage[margin=2.6cm]{geometry}
\usepackage{amsmath,amssymb,amsfonts}
\usepackage{booktabs}
\usepackage{array}
\usepackage{hyperref}
\usepackage{xcolor}
\usepackage{enumitem}

\hypersetup{
  colorlinks=true,
  linkcolor=blue!50!black,
  citecolor=blue!50!black,
  urlcolor=blue!50!black
}

\setlist[itemize]{topsep=2pt,itemsep=2pt,parsep=0pt}
\setlist[enumerate]{topsep=2pt,itemsep=2pt,parsep=0pt}
\renewcommand{\arraystretch}{1.18}

\newcommand{\TACM}{\mathrm{TACM}}
\newcommand{\Var}{\mathrm{Var}}
\newcommand{\Bin}{\mathrm{Binomial}}
\newcommand{\Unif}{\mathrm{Uniform}}
\newcommand{\Neff}{N_{\mathrm{eff}}}
\newcommand{\Qtail}{Q_{\mathrm{tail}}}
\newcommand{\KS}{\mathrm{KS}}

\title{\textbf{TACM / PRNG Diagnostics v1.3}\\
\large Interpreting Null-Calibrated Seed Sensitivity in Monte Carlo Workloads}
\author{Pascal Montmory}
\date{May 2026}

\begin{document}
\maketitle

\begin{abstract}
Version 1.2 introduced a null-calibrated diagnostic for seed-conditioned Monte Carlo variance,
\[
Z_s(f)=F_0(V_s(f)).
\]
Version 1.3 focuses on interpretation. Its goal is not to multiply indicators, nor to replace established PRNG test batteries such as TestU01, but to add a Monte Carlo workload-facing layer that quantifies and classifies seed sensitivity in a calibrated and interpretable way.

The central contribution of this revision is a three-class interpretation of calibrated diagnostics:
\[
\texttt{compatible},\qquad
\texttt{calibration\_drift},\qquad
\texttt{tail\_anomaly}.
\]
The document also adds exact binomial calibration for rare-event indicators and simple power tests based on injected defects. The intended use is practical: distinguish ordinary sampling variability, mild calibration drift, and genuine tail-sensitive anomalies in reproducible CPU-only experiments.
\end{abstract}

\section{Positioning of Versions}

\subsection{Version 1.0: Initial Framework}

Version 1.0 introduced the general structural diagnostics framework for PRNGs under Monte Carlo workloads:
\begin{itemize}
  \item seed sensitivity;
  \item functional variance;
  \item inter-seed dispersion ratios;
  \item GSQI;
  \item transition instability $\epsilon(T)$;
  \item BigCrush and cross-backend evidence;
  \item limitations and non-claims.
\end{itemize}

\subsection{Version 1.1: Development of Formulae}

Version 1.1 developed the formulae around the central empirical measure
\[
\mu_{s,b,N}
=
\frac1N
\sum_{t=1}^{N}
\delta_{u_{s,b,t}}.
\]
It clarified the role of $\chi^2$ coverage diagnostics, projected spectral gaps, effective sample size proxies, audit scores, VaR/CVaR approximations, backend reproducibility and multi-lane GPU diagnostics.

\subsection{Version 1.2: Null Calibration}

Version 1.2 introduced the calibrated statistic
\[
T_s(f)
=
\frac{(R-1)V_s(f)}
{\sigma_f^2/N},
\]
and the probability-integral-transform diagnostic
\[
Z_s(f)
=
F_{\chi^2_{R-1}}
\left(
T_s(f)
\right).
\]
Under the null model,
\[
Z_s(f)\sim \Unif[0,1].
\]
This changed the status of raw ratios. A large ratio
\[
\rho_f
=
\frac{\max_s V_s(f)}
{\min_{s,V_s(f)>0} V_s(f)}
\]
is an exploratory alert, not by itself evidence of PRNG instability.

\subsection{Version 1.3: Minimal Experimental Robustness}

Version 1.3 is an interpretation layer for $Z_s$ diagnostics. It is designed to show that TACM can distinguish:
\[
\boxed{\text{ordinary null-compatible noise}}
\]
\[
\boxed{\text{mild global calibration drift}}
\]
\[
\boxed{\text{tail anomaly}}
\]
\[
\boxed{\text{controlled injected defect}}
\]
without over-claiming from isolated p-values or raw max/min ratios.

\section{Central Message}

The central message of v1.3 is:
\[
\boxed{
\text{TACM does not replace TestU01; it adds a Monte Carlo workload-facing layer}
}
\]
\[
\boxed{
\text{that quantifies and classifies seed sensitivity in a calibrated and interpretable way.}
}
\]

TestU01-like batteries are designed to detect statistical defects in PRNG streams. TACM is aimed at a complementary question: how do seed-conditioned random streams affect actual Monte Carlo estimators under specified workloads?

\section{The Calibrated Diagnostic}

Let $s$ denote a seed, $r\in\{1,\ldots,R\}$ a replicate, $N$ the number of draws per replicate, and $f$ a diagnostic integrand. The Monte Carlo estimate is
\[
I_{s,r}(f)
=
\frac1N
\sum_{t=1}^{N}
f(u_{s,r,t}).
\]
The seed-conditioned replicated mean is
\[
\bar I_s(f)
=
\frac1R
\sum_{r=1}^{R}
I_{s,r}(f).
\]
The seed-conditioned variance is
\[
V_s(f)
=
\frac1{R-1}
\sum_{r=1}^{R}
\left(
I_{s,r}(f)-\bar I_s(f)
\right)^2.
\]

\paragraph{Important convention.}
Here $V_s(f)$ is the empirical variance of the replicated Monte Carlo estimates $I_{s,r}(f)$ under seed $s$. It is not the within-run sample variance of the raw values $f(u_t)$.

Under the idealized null model,
\[
I_{s,r}(f)
\approx
\mathcal N
\left(
\mu_f,\frac{\sigma_f^2}{N}
\right),
\]
which implies
\[
T_s(f)
=
\frac{(R-1)V_s(f)}
{\sigma_f^2/N}
\sim
\chi^2_{R-1}.
\]
The calibrated diagnostic is
\[
\boxed{
Z_s(f)
=
F_{\chi^2_{R-1}}
\left(
T_s(f)
\right).
}
\]
Under the null model,
\[
\boxed{
Z_s(f)\sim \Unif[0,1].
}
\]

\paragraph{Interpretation.}
A large raw ratio $\rho_f$ may be compatible with the null model, especially when $R$ is small. Therefore:
\[
\boxed{
\rho_f\text{ is an alert; } Z_s\text{ is the calibration.}
}
\]

\section{Interpreting the $Z_s$ Distribution}

Version 1.3 combines three readings of $Z_s$: KS tests, tail-rate diagnostics and ECDF shape.

\subsection{KS Test}

For each generator and integrand, test
\[
\{Z_s(f):s\in\mathcal B\}
\]
against $\Unif[0,1]$ and report
\[
p_{\KS}.
\]
This test is useful, but not sufficient. With many seeds, for example
\[
S=1000,
\]
the KS test may detect very small smooth deformations that are not operationally severe.

Thus
\[
p_{\KS}<0.05
\]
should not automatically be interpreted as a PRNG defect.

\subsection{Extreme Tail Rate}

Define the extreme-tail rate
\[
\Qtail(f)
=
\frac1{|\mathcal B|}
\#\{s:Z_s(f)<0.01\ \text{or}\ Z_s(f)>0.99\}.
\]
Under uniformity,
\[
\mathbb E[\Qtail(f)]\approx0.02.
\]
The interpretation is:
\[
\Qtail(f)\approx0.02
\Rightarrow
\text{no massive population of extreme seeds,}
\]
whereas
\[
\Qtail(f)\gg0.02
\Rightarrow
\text{possible tail anomaly.}
\]

\subsection{ECDF Shape}

Plot the empirical distribution function
\[
\widehat F_Z(z)
\]
against the diagonal
\[
F(z)=z.
\]
The ECDF is often more readable than histograms because it exposes smooth drift and tail deformation directly.

The v1.3 interpretation is:
\[
\boxed{
\text{ECDF close to the diagonal}\Rightarrow \texttt{compatible}.
}
\]
\[
\boxed{
\text{Smooth global ECDF deformation}\Rightarrow \texttt{calibration\_drift}.
}
\]
\[
\boxed{
\text{ECDF deformation concentrated near extremes}\Rightarrow \texttt{tail\_anomaly}.
}
\]

\section{Three-Class Classification}

Version 1.3 replaces the binary vocabulary
\[
\texttt{compatible}/\texttt{suspicious}
\]
with three operational classes:
\[
\boxed{
\texttt{compatible},\quad
\texttt{calibration\_drift},\quad
\texttt{tail\_anomaly}.
}
\]

\subsection{Compatible}

Criteria:
\[
p_{\KS}\geq0.05
\]
and
\[
\Qtail(f)\approx0.02.
\]
Interpretation:
\[
\boxed{
\text{the }Z_s\text{ values are compatible with }\Unif[0,1].
}
\]

\subsection{Calibration Drift}

Criteria:
\[
p_{\KS}<0.05
\]
but
\[
\Qtail(f)\approx0.02.
\]
Interpretation:
\[
\boxed{
\text{small global calibration deformation without excess extreme seeds.}
}
\]
This class is important because it prevents over-interpreting every KS rejection as a severe PRNG failure.

\subsection{Tail Anomaly}

Criterion:
\[
\Qtail(f)\gg0.02.
\]
Interpretation:
\[
\boxed{
\text{abnormal population of extreme seeds.}
}
\]
This is the most important class for tail-sensitive Monte Carlo workloads, risk estimation and rare-event simulation.

\section{Exact Rare-Event Calibration}

For the rare-event indicator
\[
f(u)=\mathbf 1_{\{u>0.99\}},
\]
the number of rare events in replicate $r$ under seed $s$ is
\[
K_{s,r}
=
\sum_{t=1}^{N}
\mathbf 1_{\{u_{s,r,t}>0.99\}}.
\]
Under the null model,
\[
K_{s,r}
\sim
\Bin(N,p),
\qquad
p=0.01.
\]
The normal approximation used in the $\chi^2$ variance diagnostic is often adequate for large $N$, but rare-event indicators and small samples deserve a robustness check.

Define the exact binomial calibration
\[
\boxed{
Z^{\mathrm{bin}}_{s,r}
=
F_{\Bin(N,p)}
\left(
K_{s,r}
\right).
}
\]
Because the binomial distribution is discrete, the raw CDF transform is conservative and can show ties. A randomized probability-integral transform may be used:
\[
Z^{\mathrm{bin,rand}}_{s,r}
=
F_{\Bin(N,p)}(K_{s,r}-1)
+
U_{s,r}
\cdot
\mathbb P(K=K_{s,r}),
\]
where $U_{s,r}\sim\Unif[0,1]$ is an independent calibration randomizer. This optional randomized version is closer to uniform under the exact discrete null.

The goal is not to replace the replicated-variance diagnostic, but to ensure that rare-event conclusions do not rely only on a Gaussian or $\chi^2$ approximation.

\begin{table}[h]
\centering
\caption{Rare-event exact calibration summary.}
\begin{tabular}{lrrrrl}
\toprule
Generator & $p_{\KS}^{\chi^2}$ & $p_{\KS}^{\mathrm{bin}}$ & $\Qtail$ & $N$ & Verdict \\
\midrule
MT19937 & -- & -- & -- & -- & -- \\
PCG64   & -- & -- & -- & -- & -- \\
Philox  & -- & -- & -- & -- & -- \\
SFC64   & -- & -- & -- & -- & -- \\
\bottomrule
\end{tabular}
\end{table}

\section{Experimental Validation}

\subsection{Public PRNG Baseline}

The recommended v1.3 baseline is
\[
S=1000,\qquad N=50\,000,\qquad R=8.
\]
The public PRNGs are:
\[
\text{MT19937},\quad \text{PCG64},\quad \text{Philox},\quad \text{SFC64}.
\]
The diagnostic integrands are:
\[
u,\qquad
\sin(2\pi u),\qquad
u^2,\qquad
\mathbf 1_{\{u>0.99\}}.
\]
For each generator and integrand, report:
\begin{itemize}
  \item $p_{\KS}$;
  \item $\Qtail$;
  \item $\rho_{95/5}$;
  \item the three-class verdict.
\end{itemize}

\subsection{Power Test 1: Marginal Bias}

Inject a controlled marginal perturbation, for example
\[
u'_t=(u_t+\varepsilon)\bmod1,
\]
or, for a non-measure-preserving upper-bound perturbation,
\[
u'_t=\min(1,u_t+\varepsilon).
\]
Expected diagnostics include:
\begin{itemize}
  \item $\chi^2$ marginal coverage;
  \item ECDF of $Z_s$;
  \item $\Qtail$ if the perturbation affects a tail;
  \item local drift $\Delta_U$, if available.
\end{itemize}
The goal is:
\[
\boxed{
\text{show that TACM detects a controlled marginal bias.}
}
\]

\subsection{Power Test 2: Temporal Correlation}

Inject temporal dependence:
\[
u'_t=(1-\alpha)u_t+\alpha u_{t-1}\pmod1.
\]
Expected diagnostics include:
\begin{itemize}
  \item projected $\lambda$-gap;
  \item transition instability $\epsilon(T)$;
  \item two-dimensional $\chi^2$ transition maps;
  \item possibly $Z_s$ ECDFs if the workload is affected.
\end{itemize}
The goal is:
\[
\boxed{
\text{show that TACM detects injected temporal structure.}
}
\]

\subsection{Power Test 3: Rare-Event Deletion}

For a tail-specific defect, modify the stream as follows:
\[
\text{if }u_t>0.99,\text{ replace it with probability }\alpha\text{ by }u'_t=0.98.
\]
This directly suppresses rare events above the threshold. The expected diagnostic is:
\[
\boxed{
Z^{\mathrm{bin}}\text{ detects the anomaly.}
}
\]

\section{Multiple Testing}

Across multiple generators and integrands, many tests are performed. An isolated
\[
p<0.05
\]
is therefore not enough to conclude.

For $m$ tests, a simple Bonferroni threshold is
\[
p_{\mathrm{Bonf}}=\frac{0.05}{m}.
\]
For example, with
\[
m=16,
\]
the Bonferroni threshold is
\[
p_{\mathrm{Bonf}}=0.003125.
\]
Thus:
\[
\boxed{
p_{\KS}<0.05\text{ is exploratory if multiple-testing correction does not confirm it.}
}
\]

\section{Synthesis Tables}

\begin{table}[h]
\centering
\caption{Table A: Calibration summary.}
\begin{tabular}{llrrrrrl}
\toprule
Generator & Integrand & $S$ & $N$ & $R$ & $p_{\KS}$ & $\Qtail$ & Verdict \\
\midrule
MT19937 & identity  & 1000 & 50000 & 8 & -- & -- & -- \\
MT19937 & sin       & 1000 & 50000 & 8 & -- & -- & -- \\
MT19937 & quadratic & 1000 & 50000 & 8 & -- & -- & -- \\
MT19937 & rare      & 1000 & 50000 & 8 & -- & -- & -- \\
PCG64   & identity  & 1000 & 50000 & 8 & -- & -- & -- \\
Philox  & rare      & 1000 & 50000 & 8 & -- & -- & -- \\
SFC64   & rare      & 1000 & 50000 & 8 & -- & -- & -- \\
\bottomrule
\end{tabular}
\end{table}

\begin{table}[h]
\centering
\caption{Table B: ECDF interpretation.}
\begin{tabular}{lll}
\toprule
Generator / Integrand & ECDF shape & Interpretation \\
\midrule
Smooth near diagonal & close to diagonal & \texttt{compatible} \\
Smooth shifted curve & global drift & \texttt{calibration\_drift} \\
Extreme deformation & tail excess & \texttt{tail\_anomaly} \\
\bottomrule
\end{tabular}
\end{table}

\begin{table}[h]
\centering
\caption{Table C: Power tests.}
\begin{tabular}{llll}
\toprule
Defect & Intensity & Expected diagnostic & Detected? \\
\midrule
Marginal bias & $\varepsilon$ & $\chi^2$, ECDF & yes/no \\
Temporal correlation & $\alpha$ & $\lambda$-gap, $\chi^2_{2D}$ & yes/no \\
Rare-event deletion & $\alpha$ & $Z^{\mathrm{bin}}$ & yes/no \\
\bottomrule
\end{tabular}
\end{table}

\begin{table}[h]
\centering
\caption{Table D: Rare-event exact calibration.}
\begin{tabular}{lrrl}
\toprule
Generator & $p_{\KS}^{\chi^2}$ & $p_{\KS}^{\mathrm{binomial}}$ & Comment \\
\midrule
MT19937 & -- & -- & -- \\
PCG64   & -- & -- & -- \\
Philox  & -- & -- & -- \\
SFC64   & -- & -- & -- \\
\bottomrule
\end{tabular}
\end{table}

\section{Minimal v1.3 Reproducibility Bundle}

The v1.3 bundle should extend the existing z-test pipeline with the following scripts:
\[
\texttt{plot\_z\_ecdf.py},
\quad
\texttt{compute\_binomial\_rare\_event\_z.py},
\quad
\texttt{inject\_prng\_defects.py}.
\]
Optional scripts include:
\[
\texttt{adjust\_pvalues.py},
\quad
\texttt{scan\_snr\_grid.py},
\quad
\texttt{bootstrap\_z\_scores.py}.
\]

The minimal versioned artifacts are:
\begin{verbatim}
results/z_summary.csv
figures/ecdf_Zs_identity.svg
figures/ecdf_Zs_sin.svg
figures/ecdf_Zs_quadratic.svg
figures/ecdf_Zs_rare_099.svg
figures/hist_Zs_identity.svg
figures/hist_Zs_sin.svg
figures/hist_Zs_quadratic.svg
figures/hist_Zs_rare_099.svg
\end{verbatim}
The larger files
\begin{verbatim}
data/public_prng_estimates.csv
results/z_scores_table.csv
\end{verbatim}
may be generated locally rather than versioned, depending on repository size policy.

\section{Limitations}

The v1.3 interpretation is deliberately cautious:
\begin{itemize}
  \item $Z_s$ depends on the chosen null model;
  \item the $\chi^2$ calibration assumes approximately Gaussian Monte Carlo estimates;
  \item for rare-event indicators, exact binomial or bootstrap calibration may be preferable;
  \item a non-rejected test does not prove absence of PRNG defects;
  \item public PRNGs tested in the baseline are not exhaustive;
  \item simple diagnostic integrands do not cover all industrial payoffs;
  \item injected power tests only cover controlled defect families.
\end{itemize}

\section{Conclusion}

The expected v1.3 conclusion is strong but limited:
\[
\boxed{
\text{the tested public PRNGs do not show strong seed pathologies in these simple workloads.}
}
\]
At the same time:
\[
\boxed{
\text{TACM detects, classifies and interprets raw alerts, calibration drift and tail anomalies.}
}
\]

The final message is:
\[
\boxed{
\text{TACM does not replace TestU01.}
}
\]
\[
\boxed{
\text{It adds a Monte Carlo workload-facing layer that quantifies and classifies}
}
\]
\[
\boxed{
\text{seed sensitivity in a statistically calibrated and interpretable way.}
}
\]

\end{document}